% ============================================================
%  chapters/ch26-xyloarchy-living-systems.tex
% ============================================================

\chapter{Xyloarchy and Living Systems}

\chapterepigraph{%
  A forest does not grow.\\
  It accumulates. It integrates. It forgets.\\
  It is the oldest semantic infrastructure we know.
}{Flyxion, \textit{Semantic Infrastructure}}

\sidenote{Xyloarchy:\\from Greek \textit{xylo}\\(wood) + \textit{arche}\\(governance):\\ rule by trees}

The previous chapter developed semantic infrastructure as a category of
modules and compositions. This chapter grounds that abstraction in the
most ancient and successful example of distributed knowledge infrastructure
on Earth: forest systems.

The claim is not metaphorical. Forests — through mycorrhizal networks,
chemical signaling, root competition, and structural interdependence —
implement all five knowledge evolution operations defined in Chapter~25:
extension, correction, composition, abstraction, and instantiation.
The forest is a semantic infrastructure that has been running, with
iterative refinement, for approximately 385 million years.

\section{Forest Intelligence}

The concept of plant intelligence remains controversial in biology,
partly because it is poorly defined. The accessibility framework offers
a precise version: a system exhibits intelligence to the degree that
it navigates its environment using accessibility gradients rather than
fixed reflexes.

\begin{definition}{Distributed Environmental Intelligence}{distributed-intel}
  A distributed system exhibits \emph{environmental intelligence} if
  it maintains a functional model of its environment (in the sense of
  Chapter~22) at the collective level, even when no individual component
  maintains such a model.
\end{definition}

Individual trees do not model their environment in any rich sense.
But a mycorrhizal network — the ``wood wide web'' of fungal connections
linking trees through soil — collectively integrates information about
nutrient availability, stress levels, and competitor presence across
the entire forest. The network's chemical state encodes a distributed
model of the forest environment.

\begin{historynote}
  Ibn Sina's claim that plants possess a ``vegetative soul'' ($\mathrm{nafs
  \ nabatiyya}$) — the capacity for nutrition, growth, and reproduction —
  is the earliest formal treatment of plant cognition in the philosophical
  tradition. The RSVP framework can be understood as a modern version of
  this claim, generalized: what Ibn Sina called the vegetative soul is
  the plant's accessibility navigation system, implemented through
  biochemical gradients rather than neural circuits.
\end{historynote}

\section{The Mycorrhizal Network as Semantic Module System}

The mycorrhizal network implements the operations of semantic infrastructure
(Chapter~25) in biological substrate.

\begin{itemize}
  \item \textbf{Extension:} a tree under carbon stress signals this
        through the network; neighboring trees respond by allocating
        additional photosynthate. New information (carbon stress level)
        is incorporated into the network's distributed model.
  \item \textbf{Correction:} chemical signals that indicate false alarms
        (stress signals not followed by actual damage) are modulated by
        the network's response history — a biological analog of epistemic
        updating.
  \item \textbf{Composition:} when two previously separate mycorrhizal
        networks join (as trees from different root systems become
        connected), their distributed models must be reconciled — a
        biological homotopy merge.
  \item \textbf{Abstraction:} the network's aggregate response to
        multiple individual stress signals produces a forest-level
        accessibility map — an abstraction from individual tree states.
  \item \textbf{Instantiation:} specific growth decisions (root
        elongation, leaf production, reproductive timing) instantiate
        the abstract forest-level model to individual tree behavior.
\end{itemize}

\section{Paper Mills and Epistemic Ecology}

The concept of \emph{paper mills} — organizations that mass-produce
fake academic papers for purchase — offers a striking illustration of
semantic infrastructure failure. In biological terms, paper mills are
\emph{epistemic parasites}: they exploit the nutrient channels of
academic infrastructure (journals, databases, citation networks)
without contributing to the knowledge base.

\begin{definition}{Epistemic Parasite}{epistemic-parasite}
  An \emph{epistemic parasite} is a semantic module that:
  \begin{enumerate}
    \item Occupies space in the semantic infrastructure (is cited,
          indexed, stored),
    \item Has zero or negative admissibility (its propositions are
          false or its operations are invalid),
    \item Consumes resources (attention, citation credit, storage)
          from legitimate modules.
  \end{enumerate}
\end{definition}

Mycorrhizal networks have evolved mechanisms to detect and exclude
parasitic fungi — organisms that take nutrients from the network
without contributing photosynthate. Academic institutions have analogous
mechanisms (peer review, retraction, replication), but these mechanisms
are slower and less effective at scale.

The RSVP accessibility framework predicts the consequence of epistemic
parasitism: it reduces the accessibility entropy of the knowledge base.
Each parasitic paper that is cited occupies a position in the knowledge
graph that could otherwise connect legitimate modules. The parasite does
not merely waste space — it actively reduces the connectivity of the
knowledge graph and therefore the admissible extensions available to
future researchers.

\section{Cities as Organs}

A city is an organ of civilizational metabolism. It processes inputs
(food, energy, information, people), transforms them through social and
economic processes, and exports outputs (goods, services, ideas, institutions).

\begin{definition}{City as Semantic Processor}{city-semantic}
  A city $\mathcal{U}$ is a \emph{semantic processor} if it:
  \begin{enumerate}
    \item Accumulates semantic modules (through libraries, universities,
          research institutions, cultural memory),
    \item Composes them (through the interaction of specialists from
          different domains),
    \item Abstracts from them (through the formation of professional
          traditions, schools of thought, industrial techniques),
    \item Distributes the results (through education, publication,
          trade, migration).
  \end{enumerate}
\end{definition}

The RSVP accessibility field of a city is its \emph{knowledge landscape}:
the distribution of semantic modules and their accessibility to inhabitants.
A dense, well-connected knowledge landscape (high $S_{\mathbf{Sem}}$)
allows rapid knowledge evolution; a sparse, isolated one prevents it.

Jane Jacobs' observation that cities generate innovation through the
collision of diverse activities is, in these terms, an observation about
accessibility: density and diversity increase the probability of homotopy
merges between previously separate modules.

\section{Constraint Ecology}

Living systems — forests, cities, ecosystems — are governed by
\emph{constraint ecology}: the dynamics of how constraints propagate,
compete, and evolve in complex networks.

\begin{definition}{Constraint Ecology}{constraint-ecology}
  A \emph{constraint ecology} is a dynamical system over a network of
  agents $(A_1, \ldots, A_n)$ in which:
  \begin{itemize}
    \item Each agent maintains a local constraint set $\mathcal{C}_i(t)$,
    \item Constraints propagate between connected agents through
          communication channels $E_{ij}$,
    \item The global constraint structure $\mathcal{C}(t) =
          \bigcup_i \mathcal{C}_i(t)$ evolves through
          local updates and inter-agent propagation.
  \end{itemize}
\end{definition}

Constraint ecology unifies the biological, urban, and epistemic systems
discussed in this chapter. A healthy ecosystem maintains diverse, locally
consistent constraint ecologies that can merge without catastrophic
obstruction. An unhealthy ecosystem — monoculture, epistemic bubble,
urban homogenization — has reduced constraint diversity and therefore
reduced accessibility entropy.

\begin{exercise}
  Model a simple three-tree mycorrhizal network as a constraint ecology.
  Define the state space for each tree (carbon level, water stress,
  pathogen load), the communication channels (mycorrhizal connections),
  and the constraint update rule. Show that the network reaches a more
  stable equilibrium than three isolated trees would.
\end{exercise}

\begin{exercise}
  Cities in different historical periods have had very different
  $S_{\mathbf{Sem}}$ values. Compare Baghdad at its height (9th century)
  and Baghdad under the Mongol invasion (13th century) using the
  accessibility entropy framework. What specific changes reduced
  $S_{\mathbf{Sem}}$?
\end{exercise}

\begin{exercise}[title={Open problem}]
  Define a formal measure of ``epistemic parasitism'' for academic papers:
  a score that quantifies the degree to which a paper extracts value
  from the knowledge infrastructure without contributing to it. What
  data would be needed to compute this score? Is it possible to construct
  such a measure that is not itself gameable?
\end{exercise}
